\chapter[Literature Review]{Literature Review}
\label{cp:literature}

\parindent0pt

Our system draws on six lines of work: UAVs for disaster response, consensus and its resilient variants, distributed
coverage, information spreading, the logic of knowledge, and secure localisation and safety. This chapter reviews each
briefly and ends by placing our work among them.

\section{UAV Fleets in Disaster Response}

Drones are now a standard tool in disaster management, for damage assessment, mapping, delivery of supplies and
search for victims~\citep{Erdelj2017}. The common architectures are either centralised, where a ground station plans
for all drones, or decentralised, where each drone plans locally from shared information. Centralised planning breaks
down exactly when it is needed most: with ground communication lost, and with drones often out of range of one
another. We therefore follow the decentralised model, in which every drone acts on its own sensors and on information
relayed over a time-varying, multi-hop radio graph.

\section{Consensus and Resilient Consensus}
\label{sec:lit-consensus}

\textbf{Consensus.} \citet{OlfatiSaber2007} survey consensus protocols in which each agent repeatedly moves its state
towards the states of its neighbours. On a connected graph, every agent converges to a common value, and the rate is
governed by the algebraic connectivity $\lambda_2$ of the graph Laplacian~\citep{Fiedler1973, Mesbahi2010}. Formation
control, rendezvous and distributed estimation are all built on this primitive~\citep{Bullo2009}.

\textbf{Byzantine faults.} \citet{Pease1980} and \citet{Lamport1982} showed that agreement in the presence of $F$
arbitrarily misbehaving (Byzantine) agents requires $N > 3F$ agents and enough connectivity, and that signed messages
make the problem much easier. Their model, in which faulty agents can send arbitrary but authenticated messages, is
the one we adopt for the radio layer.

\textbf{W-MSR and network robustness.} \citet{LeBlanc2013} introduced the W-MSR algorithm. Each normal agent discards up
to $F$ neighbour values above its own and up to $F$ below, then averages the rest. They proved that normal agents reach
consensus inside the convex hull of their initial values if and only if the graph is sufficiently \emph{robust}. They
introduced $r$-robustness and $(r,s)$-robustness as the right connectivity notions, replacing the classical vertex
connectivity. \citet{Saldana2017} extended the guarantee to time-varying networks of mobile robots.
\citet{Usevitch2020} gave mixed-integer programs to compute robustness, which is coNP-hard in general.
\citet{Mitra2019} designed Byzantine-resilient distributed observers for estimating a dynamic state. W-MSR is our
\naive{} baseline: it is exactly the resilience a well-designed fleet from this literature would have.

\section{Distributed Coverage and Task Allocation}

\citet{Cortes2004} formulated coverage control as moving agents to the centroids of their Voronoi cells, a rule that
locally optimises a coverage cost and needs only neighbour information. Later work replaced Euclidean distance with
geodesic distance in non-convex environments, so that a cell behind a wall belongs to the agent that can actually
reach it. Our fleet uses a \emph{geodesic Voronoi partition} computed by multi-source Dijkstra~\citep{Dijkstra1959}
over the drone's own map. Each drone computes it locally, and work-stealing balances the load when a region empties.

\section{Information Spreading}

Gossip, or epidemic, protocols spread information through a network by having each node forward new items to its
neighbours~\citep{Demers1987}. Their simple-contagion dynamics follow the SI model taught in the course. We use gossip for
all fleet-wide news (cleared cells, survivors, terrain, accusations). We show that accusations of lying cannot spread
as a simple contagion. Because a single accuser could itself be the liar, belief in an accusation must follow a
\emph{threshold} rule (\autoref{sec:accusations}).

\section{Knowledge and Possible Worlds}
\label{sec:lit-epistemic}

Epistemic logic formalises what an agent \emph{knows}. In Kripke semantics~\citep{Kripke1963}, an agent knows $\varphi$
in a world $w$ if $\varphi$ is true in every world the agent cannot distinguish from $w$. When indistinguishability is an
equivalence relation, this is the logic S5. \citet{Fagin1995} developed this theory for distributed systems.
\citet{Halpern1990} used it to prove that common knowledge cannot be attained over unreliable channels.
\citet{vanDitmarsch2007} treat how knowledge changes as evidence arrives, which is our setting: every proof of lying is
a public or private announcement that removes worlds. Knowledge-based programs, which condition actions on knowledge
tests such as ``if $\Kn_i\varphi$ then act'', are a central idea in that literature. Our motion gate is a direct
instance of one.

\section{Secure Localisation, Spoofing and Safety}

\textbf{Verifying positions.} A drone that lies about its position can be caught only with a measurement it cannot
forge. \citet{Capkun2005} showed that time-of-flight ranging (as in ultra-wideband radios) gives \emph{verifiable}
distance bounds, because a node can make itself appear farther away by delaying its reply, but never closer. With
three or more verifiers around it, a position claim can be checked by multilateration. \citet{Gil2017} used physical
properties of wireless signals to detect spoofed robot identities in multi-robot networks. Our proof rules R1 (range)
and R6 (shared range) follow the same principle: a claim is checked against a physical measurement taken by the
listener.

\textbf{Collision avoidance.} Velocity obstacles~\citep{Fiorini1998} and control barrier functions~\citep{Ames2019}
are standard tools for keeping robots apart. Both assume that the robot knows where its neighbours are. When a
neighbour may be lying, its reported position cannot be trusted. Our knowledge gate therefore checks a manoeuvre
against \emph{every} position a neighbour could occupy in \emph{every} world the drone cannot rule out.

\section{Runtime Verification and Sequential Testing}

Signal Temporal Logic (STL)~\citep{Maler2004} expresses requirements over continuous signals, such as ``always keep
1\,m separation'' or ``eventually find survivor 7''. Its \emph{robustness} semantics~\citep{Donze2010} returns a real
number whose sign says whether the requirement held and whose size says by how much. We use STL robustness to grade
every run offline. For deciding whether a found survivor is alive, we use Wald's Sequential Probability Ratio
Test~\citep{Wald1945}. It needs the fewest observations on average for given error rates, so a drone hovers only as
long as needed.

\section{Positioning of this Work}

\begin{table}[h]
\centering
\renewcommand{\arraystretch}{1.25}
\begin{tabularx}{\linewidth}{@{}X c c c c@{}}
\toprule
\textbf{Approach} & \textbf{Numeric} & \textbf{Claims} & \textbf{Identifies} & \textbf{Safe w.r.t.} \\
                  & \textbf{resilience} & \textbf{verified} & \textbf{liars} & \textbf{lying peers} \\
\midrule
Plain consensus~\citep{OlfatiSaber2007}           & \xmark & \xmark & \xmark & \xmark \\
W-MSR / robust networks~\citep{LeBlanc2013}       & \cmark & \xmark & \xmark & \xmark \\
Spoof detection~\citep{Gil2017}                    & \xmark & partly & \cmark & \xmark \\
Collision avoidance~\citep{Fiorini1998,Ames2019}   & \xmark & \xmark & \xmark & \xmark \\
\textbf{VeriFly gated fleet (ours)}                & \cmark & \cmark & \cmark & \cmark \\
\bottomrule
\end{tabularx}
\caption{What each line of work provides for a search fleet with lying members.}
\label{tab:positioning}
\end{table}

To our knowledge, no prior system combines resilient numeric consensus with exact possible-worlds knowledge of
\emph{who} is lying, provenance-tagged shared maps, and a knowledge-based motion gate, all evaluated in a decentralised
3D search mission with a controlled, paired comparison. The contribution is not a new consensus algorithm. It is a
demonstration, with numbers, that the decision layer above consensus is where a search fleet's resilience is decided.
